%%%%%%%%%%%%%%%%%
% CV tailored for ProjectTech Engineering - Technical Data & Design Quality Process Manager
%%%%%%%%%%%%%%%%%

\documentclass[8pt,a4paper,ragged2e,withhyper]{altacv}

\geometry{left=1.25cm,right=1.25cm,top=.5cm,bottom=1.cm,columnsep=1.2cm}

\usepackage{paracol}
\usepackage{fontawesome5}
\usepackage{hyperref}

\iftutex 
  \setmainfont{Roboto Slab}
  \setsansfont{Lato}
  \renewcommand{\familydefault}{\sfdefault}
\else
  \usepackage[rm]{roboto}
  \usepackage[defaultsans]{lato}
  \renewcommand{\familydefault}{\sfdefault}
\fi

\definecolor{SlateGrey}{HTML}{2E2E2E}
\definecolor{LightGrey}{HTML}{666666}
\definecolor{DarkPastelRed}{HTML}{8AC0D9}
\definecolor{PastelRed}{HTML}{00B0DA}
\definecolor{GoldenEarth}{HTML}{B4AD0C}

\colorlet{name}{black}
\colorlet{tagline}{PastelRed}
\colorlet{heading}{DarkPastelRed}
\colorlet{headingrule}{GoldenEarth}
\colorlet{subheading}{PastelRed}
\colorlet{accent}{PastelRed}
\colorlet{emphasis}{SlateGrey}
\colorlet{body}{LightGrey}

\definecolor{violet}{HTML}{5A2582}
\definecolor{rouge}{HTML}{F53B3B}
\definecolor{noir}{HTML}{00232C}
\definecolor{bleu}{HTML}{00B0DA}
\definecolor{rose}{HTML}{F831A0}
\definecolor{jaune}{HTML}{FFEC00}
\definecolor{vert}{HTML}{34AD6C}

\renewcommand{\namefont}{\Huge\rmfamily\bfseries}
\renewcommand{\personalinfofont}{\footnotesize}
\renewcommand{\cvsectionfont}{\LARGE\rmfamily\bfseries}
\renewcommand{\cvsubsectionfont}{\large\bfseries}
\renewcommand{\cvItemMarker}{{\small\textbullet}}
\renewcommand{\cvRatingMarker}{\faCircle}

\input{pubs-num.tex}
\addbibresource{sample.bib}

\begin{document}

\name{BOILEAU Guillaume}
\tagline{Technical Data \& Design Quality Process Manager | Traceability, Data Governance, Validation \& Process Structuring}

\photoL{2.8cm}{moi}

\personalinfo{%
  \email{guillaume.boileaupro@gmail.com}
  \phone{+33 6 59 94 85 84}
  \location{Alpes-Maritimes, FRANCE}
  \linkedin{guillaume-boileau-phd-891b81265}
  \researchgate{Guillaume-Boileau-2}
  \orcid{0000-0002-3576-6968}
}

\makecvheader
\vspace{-0.3cm}

\AtBeginEnvironment{itemize}{\small}
\columnratio{0.64}

\begin{paracol}{2}

\cvsection{Experience}

\cvevent{\textbf{Software Engineer - System Integration \& Operational Data Interfaces}}{VEV Platform Services France}{2025 -- 2026}{Valbonne, France}

\textbf{Context:} Development and integration of software services for EV charging infrastructure within a CPMS platform connected to operational hardware systems.

\textbf{Objective:} Support reliable operational data flows, software/hardware interface consistency and service behaviour in a technology-driven industrial environment.

\begin{itemize}
  \item \textbf{Technical data flows:} Implemented and monitored FTP-based operational data exchange services connected to EV charging infrastructure.
  \item \textbf{Data consistency:} Checked consistency between platform services, operational data flows and field equipment behaviour.
  \item \textbf{System integration:} Contributed to the integration and maintenance of software services interacting with charging station systems.
  \item \textbf{Hardware/software interfaces:} Analysed interactions between backend services, operational data repositories and connected infrastructure.
  \item \textbf{Operational validation:} Verified service behaviour, data exchanges and interface continuity in production-like conditions.
  \item \textbf{Incident analysis:} Investigated production issues involving software behaviour, data consistency and hardware-connected systems.
  \item \textbf{Software development:} Developed web and backend services using TypeScript, Angular and Node.js.
  \item \textbf{Agile tooling:} Used Zenhub for task tracking, backlog follow-up, iterative development and technical coordination.
\end{itemize}

\divider

\cvevent{\textbf{Senior R\&D Engineer - Technical Data, Validation \& Process Structuring}}{Laboratoire Lagrange, Observatoire de la Côte d'Azur, CNES}{2023 -- 2025}{Nice, France}

\textbf{Context:} Engineering activities for the LISA ESA/NASA space mission, involving complex software chains, model integration, simulation workflows, validation campaigns and international coordination.

\textbf{Objective:} Structure, integrate and validate technical data, models and software workflows under strong consistency, traceability, reproducibility and verification constraints.

\begin{itemize}
  \item \textbf{Technical data structuring:} Structured complex technical datasets, simulation outputs, model parameters and result repositories for mission-level analyses.
  \item \textbf{Traceability:} Maintained traceability between input assumptions, software versions, model configurations, generated outputs and validation conclusions.
  \item \textbf{Data governance logic:} Defined practical rules for data organisation, naming, versioning, reproducibility and comparison across multiple contributors.
  \item \textbf{Design quality logic:} Supported quality-oriented validation of models, workflows and deliverables through consistency checks and documented review processes.
  \item \textbf{Validation workflows:} Built integration, verification and validation workflows for complex models, software chains and simulation outputs.
  \item \textbf{Process formalisation:} Formalised comparison campaigns, acceptance logic, technical checks and validation evidence for collaborative engineering activities.
  \item \textbf{Software platform:} Developed CATpyLISA, a Python platform for large-scale model integration, comparison, validation and technical review.
  \textcolor{DarkPastelRed}{\href{https://gitlab.oca.eu/gboileau/lisa_l3_tools}{\bf GitLab:CATpyLISA}}
  \item \textbf{Documentation management:} Used Confluence to support technical documentation, validation follow-up, coordination notes and knowledge sharing.
  \item \textbf{Cross-functional coordination:} Coordinated contributors, supervised interns and aligned technical activities across scientific, software and validation stakeholders.
  \item \textbf{Continuous improvement:} Improved workflows, documentation practices and analysis reliability to reduce inconsistencies and strengthen reproducibility.
\end{itemize}

\divider

\cvevent{\textbf{Data Scientist - Diagnostics, Data Quality \& Performance Validation}}{Universiteit Antwerpen, Belgium}{2021 -- 2023}{Antwerp, Belgium}

\textbf{Context:} Data analysis and performance studies for precision instruments in a high-requirement technical environment linked to gravitational-wave detector performance.

\textbf{Objective:} Identify, characterize and mitigate sources affecting system sensitivity using robust statistical, computational and validation-oriented methods.

\begin{itemize}
  \item \textbf{Data quality analysis:} Analysed measurement and simulation outputs to detect abnormal behaviours, inconsistencies and degradation sources.
  \item \textbf{Technical diagnostics:} Identified contributors to performance degradation and supported prioritisation of mitigation actions.
  \item \textbf{Performance validation:} Assessed sources affecting system sensitivity, measurement quality and performance stability.
  \item \textbf{Analysis pipelines:} Developed Python-based pipelines for noise source identification, characterization and mitigation.
  \item \textbf{Traceable reporting:} Produced reproducible and documented technical analyses for international engineering and research stakeholders.
  \item \textbf{System impact:} Contributed to studies identifying environmental and instrumental constraints for next-generation gravitational-wave observatories.
\end{itemize}

\divider

\cvevent{\textbf{Junior R\&D Engineer - Signal Analysis, Simulation \& Software Validation}}{ARTEMIS, Observatoire de la Côte d'Azur}{2018 -- 2021}{Nice, France}

\textbf{Context:} Engineering and analysis activities for gravitational-wave mission preparation, involving signal analysis, simulation, software workflows and noise characterization.

\textbf{Objective:} Develop robust analysis methods for complex signals and support technical investigations in demanding system environments.

\begin{itemize}
  \item \textbf{Simulation data:} Built simulation approaches for complex signal environments, uncertainty studies and large-scale datasets.
  \item \textbf{Software validation:} Compared simulated outputs, model behaviours and analysis results to assess consistency and robustness.
  \item \textbf{Technical investigations:} Investigated noise sources through cross-sensor correlation analyses in diagnostic contexts.
  \item \textbf{Signal analysis:} Developed methods for the separation of overlapping low-SNR signals in complex datasets.
  \item \textbf{Mission performance:} Supported the identification of limiting factors affecting mission-level performance and system sensitivity.
\end{itemize}

\switchcolumn

\cvsection{Profile for ProjectTech}

\textbf{Engineer with strong experience in technical data structuring, validation workflows, documentation control, traceability and cross-functional coordination in complex technology-driven environments}. Background developed through space systems, industrial software interfaces and large-scale collaborative engineering activities.

For a \textbf{Technical Data \& Design Quality Process Manager} role, I bring a practical engineering mindset focused on securing the \textbf{reliability, consistency and traceability of technical data}, structuring validation processes, formalising documentation practices and supporting technical teams across R\&D, software, operations and system-level activities.

\begin{itemize}
  \item Strong experience structuring complex technical data, simulation outputs, model repositories and operational data flows.
  \item Ability to ensure consistency between design assumptions, software outputs, documentation and operational behaviour.
  \item Experience defining traceability rules across inputs, versions, configurations, outputs, validation evidence and technical conclusions.
  \item Quality-oriented engineering mindset aligned with ISO 9001 principles: documentation, version control, traceability, reproducibility, continuous improvement and process formalisation.
  \item Practical experience with validation workflows, technical reviews, anomaly analysis, data consistency checks and documented acceptance logic.
  \item Cross-functional coordination with R\&D, software, engineering, operations and international technical stakeholders.
  \item Ability to support technical teams in improving data reliability, documentation practices and process discipline.
  \item Strong analytical mindset, structured reporting and technical communication in English.
\end{itemize}

\cvsection{Languages}

\cvskill{English}{4}
\divider
\cvskill{Spanish}{2}
\divider

\medskip

\cvsection{Education}

\cvevent{PhD in Physics}{Université Côte d'Azur}{2018 -- 2021}{Nice, France}

\divider

\cvevent{Selective Graduate Program in Fundamental Physics (Magistère)}{Université Paris-Saclay}{2016 -- 2018}{Orsay, France}

\divider

\cvevent{MSc in Astronomy and Astrophysics}{Observatoire de Paris / Université Paris-Saclay}{2017 -- 2018}{Paris, France}

\divider

\cvevent{BSc in Fundamental Physics}{Université Paris-Saclay}{2015 -- 2016}{Orsay, France}

\cvsection{Professional Training}

\cvevent{Supervised Machine Learning (Regression \& Classification)}{DeepLearning.AI - Andrew Ng}{2020}{Coursera}

\end{paracol}

\cvsection{Projects}

\cvevent{CATpyLISA - Technical Data, Model Integration \& Validation Platform}{Python Platform for the LISA ESA/NASA Space Mission}{2023 -- 2025}{}

Development of a Python platform dedicated to large-scale model integration, comparison, validation and technical review for LISA mission-level data analysis workflows.

\textbf{Role:} Technical data structuring, model integration, software workflow design, validation campaigns, documentation support, consistency checks and coordination with multiple contributors.

\textbf{Relevance for technical data and design quality:} Work involving structured repositories, versioned data products, traceability between configurations and outputs, documented validation evidence, anomaly investigation and continuous improvement of collaborative workflows.

\divider

\cvevent{Co-author and full member of the LISA Consortium}{ESA/NASA Space Mission - Large-scale International Collaboration}{2018 -- now}{}

Contribution to software, analysis and validation workflows for the LISA space mission within an international engineering and scientific collaboration involving ESA, NASA and multiple research institutes.

\textbf{Role:} Software workflows, technical coordination, model integration, simulation campaigns, validation activities, complex signal-processing tasks and collaborative engineering support.

\textbf{Program scale:} Major international space mission with an estimated budget of approximately 2 billion EUR over 10 years.

\divider

\cvevent{Co-author and full member of Einstein Telescope and LIGO--Virgo--KAGRA Collaborations}{International Collaborations}{2021 -- now}{}

Contribution to collaborative developments in data analysis, signal characterization, software methods and technical investigations within the LIGO--Virgo--KAGRA and Einstein Telescope collaborations.

\textbf{Role:} Development of analysis tools, contribution to technical activities and support to complex software workflows in high-standard collaborative environments.

\textbf{System impact:} Studies on environmental and instrumental noise sources helping identify fundamental design constraints and performance limitations for next-generation gravitational-wave observatories.

\cvsection{Strengths}

\vspace{-.3cm}

\cvsubsection{Technical Data, Traceability \& Quality Processes}

{\small
\cvtag{Technical Data Management}
\cvtag{Data Governance}
\cvtag{Data Structuring}
\cvtag{Data Repositories}
\cvtag{Data Consistency}
\cvtag{Technical Databases}
\cvtag{Configuration Control}
\cvtag{Version Control}
\cvtag{Traceability}
\cvtag{Documentation Control}
\cvtag{Archiving Logic}
\cvtag{Validation Evidence}
\cvtag{Quality Processes}
\cvtag{ISO 9001 Logic}
\cvtag{Continuous Improvement}
\cvtag{Process Formalisation}
}

\divider\smallskip
\vspace{-.3cm}

\cvsubsection{Design-to-Production, Validation \& Interfaces}

{\small
\cvtag{R\&D Interface}
\cvtag{Industrial Interfaces}
\cvtag{Design Quality}
\cvtag{Design-to-Production Logic}
\cvtag{System Integration}
\cvtag{Software Integration}
\cvtag{Hardware / Software Interfaces}
\cvtag{Operational Data Flows}
\cvtag{Validation Workflows}
\cvtag{Acceptance Criteria}
\cvtag{Test Data Analysis}
\cvtag{Anomaly Analysis}
\cvtag{Root Cause Analysis}
\cvtag{Performance Indicators}
\cvtag{Technical Diagnostics}
\cvtag{Operational Reliability}
}

\divider\smallskip
\vspace{-.3cm}

\cvsubsection{Tools, Software \& Coordination}

{\small
\cvtag{Python}
\cvtag{Linux}
\cvtag{Bash}
\cvtag{Git}
\cvtag{GitLab}
\cvtag{GitHub}
\cvtag{Docker}
\cvtag{CI/CD}
\cvtag{SQL}
\cvtag{TypeScript}
\cvtag{Angular}
\cvtag{Node.js}
\cvtag{REST API}
\cvtag{FTP}
\cvtag{Confluence}
\cvtag{Zenhub}
\cvtag{Agile Tooling}
\cvtag{Technical Documentation}
\cvtag{Technical Coordination}
\cvtag{Stakeholder Communication}
\cvtag{Technical English}
}

\end{document}